% !TEX root = ../../main_docker.tex
% Dateisystem, Volumes, Images

\section{Dateisystem, Volumes, Repositories}
\label{sec:fs}

\subsection{Das geschichtete Container-Dateisystem}
\label{sec:overlay}
Das virtuelle Dateisystem eines Containers besteht aus übereinander
gelegten Schichten:
\begin{ztab}{@{}lX@{}}
oben & \T{Container-Layer}, read-write -- die Änderungen \\
Mitte & weitere \emph{Image}-Schichten: nachinstallierte Software,
  Webserver, DB, dotnet \dots{} \\
unten & \T{Base-Image}, meist ein minimales Linux (Ubuntu, Debian) \\
\end{ztab}
Alle Schichten ausser der obersten sind \T{read-only}.
\T{Jeder Container bekommt seinen eigenen Container-Layer} -- auch mehrere
aus demselben Image. Der Layer wird \emph{mit} dem Container gelöscht.
\merkzeile Praktisch für reinen Tisch (temporäre Dateien verschwinden),
unpraktisch für Software, die Daten dauerhaft speichern muss.

\SW{Wo das physisch liegt} \R{/var/lib/docker/overlay2/\dots} und
\R{/var/lib/containerd/\dots}. Suchen mit
\R{sudo find /var/lib/ -name <datei>}.

\begin{falle}[Zwei Beobachtungen dazu]
1.~Bei einem \emph{laufenden} Container findet \R{find} die Datei
\T{zweimal} (docker und containerd), bei einem gestoppten nur einmal.\\
2.~Diese Verzeichnisse werden von Docker verwaltet und dürfen
\textbf{nie} direkt verändert werden.
\end{falle}

\subsection{Volumes}
\label{rez:volume}
Ein Volume bildet einen Pfad des Host-Dateisystems in das
Container-Dateisystem ab. Änderungen sind auf beiden Seiten sofort
sichtbar, und die Dateien \T{überleben das Löschen des Containers} --
alles andere im Container-Layer nicht.

\begin{ztab}{@{}lX@{}}
\T{Host Volume} & \R{-v /pfad/im/host:/pfad/im/cnt} -- Nutzer bestimmt
  beide Seiten \\
\T{Anonymes V.} & \R{-v /pfad/im/cnt} -- Docker legt das Host-Verzeichnis
  unter \R{/var/lib/docker/volumes/<hash>/\_data/} an; kaum gebraucht \\
\T{Named Volume} & \R{-v name:/pfad/im/cnt} -- wie anonym, aber mit Namen \\
\end{ztab}

\begin{bsp}[Wofür Volumes wirklich da sind]
Versions-Update eines Datenbankservers: Die containerisierte DB legt ihre
Daten im Volume ab. Beim Update wird der alte Container gestoppt und
gelöscht -- die Daten bleiben. Der neue Container aus dem aktuellen Image
greift auf dieselben Daten zu.
\end{bsp}

\SW{Administration} \R{docker volume ls} $\cdot$
\R{docker volume rm <vol>} $\cdot$ \R{docker volume prune}.
Volumes überleben ihren Container, müssen also von Hand aufgeräumt werden.

\begin{falle}[Warum \R{volume prune} trügerisch ist]
Es löscht \emph{alle} Volumes, die gerade keinem Container zugeordnet
sind. Genau die wertvollen Daten eines Volumes, dessen Container du
versehentlich schon gelöscht hast, sind damit weg -- und es gibt
\textbf{kein} Docker-Kommando, um vorher den Inhalt eines Volumes
anzusehen.
\end{falle}

\subsection{Image aus Container: \R{docker commit}}
\begin{lstlisting}
docker commit <container> <image>
\end{lstlisting}
Friert die Änderungen eines laufenden oder gestoppten Containers als neue
Image-Schicht ein. Damit lässt sich manuell installierte Software in ein
Image übernehmen: Ubuntu starten $\to$ \R{apt install nano} $\to$
\R{commit} $\to$ das neue Image enthält nano.
\merkzeile Das ist die Antwort auf „Kann aus einem Container ein Image
erzeugt werden?“ -- ja, mit \R{commit}. Der saubere Weg bleibt aber das
Dockerfile.

\subsection{Repositories und Docker Hub}
\label{rez:tag}
\T{Docker Hub} ist das Standard-Repository: \R{docker run} und
\R{docker pull} holen Images von dort. Öffentliche Repositories sind für
jeden zugänglich, private nur für den Besitzer (bei Docker Hub ab einer
gewissen Anzahl kostenpflichtig) -- oder on premise mit eigener
Repository-Software (CNCF Distribution).

\begin{ztab}{@{}lX@{}}
\R{docker search <begriff>} & Images suchen; Spalten \R{STARS},
  \R{OFFICIAL}, \R{AUTOMATED}. Offizielle Images sind vertrauenswürdiger. \\
\R{docker pull <img>} & herunterladen, ohne zu starten \\
\end{ztab}

\begin{formel}[Aufbau eines Image-Namens]
\begin{center}\ttfamily
<repo>/<source>/<image>:<tag>
\end{center}
\begin{tabularx}{\linewidth}{@{}lX@{}}
\R{<image>} & der Name, z.\,B. \R{ubuntu}, \R{hello-world} \\
\R{<source>} & Nutzer oder Organisation auf Docker Hub;
  Standard \R{library} (die offiziellen Images) \\
\R{<tag>} & die Version; Standard \R{latest} \\
\R{<repo>} & die Registry; Standard \R{index.docker.io} \\
\end{tabularx}\\[2pt]
Also ist \R{docker run hello-world} identisch mit\\
\R{docker run index.docker.io/library/hello-world:latest}.\\
Weitere Beispiele: \R{ubuntu:jammy} (22.04),
\R{kaohslu/epp-demo}, \R{mcr.microsoft.com/dotnet/runtime:5.0}.
\end{formel}

\begin{falle}[\R{run} aktualisiert nicht]
\R{docker run <img>} nimmt die \T{neueste lokal vorhandene} Version --
auch wenn im Repository längst eine neuere liegt. Nur \R{docker pull}
holt die aktuelle.\\
Und beim \R{pull}: Das alte Image wird \emph{nicht} überschrieben,
sondern verliert nur seinen Tag und erscheint mit \R{<none>}
(\emph{dangling image}) -- Container könnten es noch brauchen.
\end{falle}

\subsection{Der Weg eines eigenen Images}
\begin{lstlisting}
docker build -t <user>/<name> .   # Image bauen und taggen
docker login -u <user>            # an Docker Hub anmelden
docker image push <user>/<name>   # hochladen
docker run <user>/<name>          # auf beliebigem Host starten
\end{lstlisting}
Drei Zeilen zum Veröffentlichen, eine zum Benutzen -- ohne dass auf dem
Zielrechner irgendeine Laufzeitumgebung installiert sein muss.
